\section{Objetivos}

O objetivo principal deste trabalho é avaliar o problema de otimização de plano de voo de veículos aéreos não tripulados (UAVs) com foco na coleta de dados de sensores IoT distribuídos em uma área geográfica a partir de uma abordagem matemática.

\subsection*{Objetivos específicos}

\begin{itemize}
    \item Formular o problema utilizando tempo discretizado e variáveis de decisão apropriadas.
    \item Modelar múltiplos objetivos, como a  a eficiência energética (dados coletados por unidade de energia), idade da informação (AoI) dos sensores, e tempo total de voo do UAV.
%    \begin{itemize}
%        \item Maximização da eficiência energética (dados coletados por unidade de energia).
%        \item Minimização da idade da informação (AoI) dos sensores.
%        \item Minimização do tempo total de voo do UAV.
%    \end{itemize}
%    \item Incorporar restrições operacionais relevantes, como:
%    \begin{itemize}
%        \item Limites de energia disponíveis.
%        \item Restrições de comunicação (raio de alcance dos sensores).
%        \item Restrições de deslocamento baseadas em velocidade máxima.
%    \end{itemize}
%    \item Comparar diferentes abordagens para formulação multiobjetivo, incluindo soma ponderada e método $\varepsilon$-restrito.
    \item Implementar a formulação em um resolvedor de otimização (como Gurobi) e realizar experimentos computacionais, comparando resultados ao que já existe na literatura.
\end{itemize}
